\begin{table}[htbp]
\centering
\small
\caption{Combinations deliberately excluded from the CSCAS baseline grid.}
\label{tab:cscas-grid-exclusions}
\begin{tabular}{p{0.32\textwidth}p{0.58\textwidth}}
\toprule
\textbf{Excluded combination} & \textbf{Reason} \\
\midrule
Zero-shot $\times$ full test set &
Would require $\sim$1.26M individual LLM prompts; not computationally
tractable. Kept on the shared 20k subsample like every other
non-replication baseline. \\[4pt]
Text models (BERT / SecureBERT / zero-shot) $\times$ full (42-feature) schema &
The 34 \texttt{*Similarity} columns are non-textual numeric scores;
serializing them into the input string was judged to add token cost
without giving a transformer usable signal (subword tokenization does not
expose numeric structure). Text models keep the 6-field schema and are
swept only on the eval-set axis. \\[4pt]
BERT-embedding + similarity-vector hybrid architecture &
Would let a text model genuinely use the similarity features, but
constitutes a new model architecture rather than a variant of the existing
baseline -- scoped out as future work rather than added to this grid. \\[4pt]
Classifiers (RF, LogReg, XGBoost, mining variants) $\times$ full$-$SCAS /
full$+$SCAS schema &
Those two variants exist specifically to test whether handing a
\emph{precomputed outlier flag} to a one-class \emph{anomaly} detector is
circular. Not a meaningful ablation for supervised classifiers, which
already train on the paper's unmodified 42-feature \texttt{full} schema
(SCAS included, matching the paper's own protocol). \\[4pt]
\texttt{SignatureID} as a feature, anywhere &
Nominal row identifier, not a real signal -- excluded from every schema
(base, full, and both anomaly full variants) for every model, including
the paper-replication RF, where it is dropped from \texttt{base} but
retained in \texttt{full} per the paper's own protocol. \\
\bottomrule
\end{tabular}
\end{table}
